\section{Conclusiones}

El desarrollo de este Trabajo de Fin de Grado ha culminado en una aplicación web funcional para el análisis y la simulación educativa de inversiones, desplegada en entorno real y validada tanto desde el punto de vista técnico como visual. El resultado final no se limita a una demostración conceptual, sino que incorpora autenticación, persistencia, varios modos de uso y una experiencia de interfaz unificada tras un rediseño completo por fases.

Entre las principales conclusiones del proyecto pueden destacarse las siguientes:

\begin{itemize}
    \item \textbf{Cumplimiento de objetivos funcionales y técnicos:} se ha desarrollado una herramienta capaz de combinar análisis histórico, simulación progresiva y apoyo pedagógico, pero además con una implementación real basada en Flask, Jinja2, Peewee, PostgreSQL, Render, GitHub Actions y consumo de datos históricos a través de yfinance. El sistema integra un modo práctica, un modo carrera, un modo horizonte experimental y un Tutor IA opcional asociado al informe final.

    \item \textbf{Evolución desde prototipo a sistema con persistencia y continuidad:} el proyecto comenzó con una base funcional más simple y terminó incorporando autenticación por correo, distinción entre usuario invitado y registrado, historial persistente, almacenamiento de sesiones de carrera, reanudación de partidas y separación entre sesión Flask y persistencia relacional.

    \item \textbf{La complejidad principal del trabajo se concentró en la integración de flujos internos:} más allá de la interfaz, el valor técnico del proyecto reside en haber articulado validación de entradas, descarga de mercado, control de errores, simulación por turnos, benchmark, eventos pedagógicos, persistencia híbrida y generación de informes finales dentro de una única aplicación web coherente.

    \item \textbf{El Tutor IA aportó valor pedagógico, pero también exigió restricciones técnicas y éticas explícitas:} su implementación obligó a diseñar comprobación de disponibilidad por entorno, construcción de payload resumido, prompt acotado, timeout, tratamiento de errores, salida JSON y renderizado seguro en frontend. Esto demuestra que su integración no consistió únicamente en invocar una API externa, sino en encapsularla dentro de un flujo controlado y prudente.

    \item \textbf{El producto final ha quedado razonablemente validado, aunque no exhaustivamente verificado en todos los planos:} la aplicación quedó respaldada por validaciones locales repetibles, integración continua mediante GitHub Actions, despliegue estable en Render y una batería final de \textbf{29 pruebas superadas}. Sin embargo, también persisten limitaciones reconocidas, como la ausencia de pruebas end-to-end completas y la dependencia de proveedores externos como Yahoo Finance y OpenAI.

    \item \textbf{La trazabilidad del uso de IA debe entenderse como parte del propio objeto de evaluación:} el historial Git visible refleja un uso intensivo de herramientas asistidas en la fase final del proyecto. Lejos de ocultarlo, la conclusión razonable es que el valor académico del trabajo debe juzgarse por la capacidad del autor para definir objetivos, integrar cambios, validar resultados, documentar la implementación y defender con solvencia el funcionamiento interno del sistema entregado.
\end{itemize}

En conjunto, el proyecto demuestra que es posible construir, dentro del marco de un TFG, una aplicación web con valor pedagógico real, apoyo en datos históricos, despliegue operativo y un grado de madurez superior al de un prototipo meramente académico.

\section{Trabajo futuro}

Aunque el sistema final alcanzado cubre de forma satisfactoria los objetivos principales del proyecto, siguen existiendo líneas de mejora realistas y coherentes para futuras iteraciones:

\subsection{Mejoras técnicas}
\begin{itemize}
    \item \textbf{Pruebas end-to-end:} incorporar una capa adicional de pruebas E2E con herramientas como Playwright o equivalentes para automatizar escenarios completos de navegación, autenticación y validación visual.
    \item \textbf{Accesibilidad avanzada:} profundizar en contraste, navegación por teclado, ayudas semánticas y revisión sistemática de accesibilidad para asegurar una experiencia más inclusiva.
    \item \textbf{Monitorización y observabilidad:} añadir mecanismos de monitorización del despliegue, alertas y trazas que faciliten el diagnóstico en producción.
    \item \textbf{Despliegue sobre infraestructura no gratuita:} migrar a una solución con menos restricciones de expiración, memoria o disponibilidad para reducir la fragilidad asociada a los planes gratuitos.
    \item \textbf{Seguridad y gestión de usuarios:} ampliar las capacidades de recuperación de cuenta, endurecimiento de sesiones y gestión de credenciales si el proyecto evolucionara hacia un producto más abierto al público.
\end{itemize}

\subsection{Mejoras funcionales y pedagógicas}
\begin{itemize}
    \item \textbf{Más contenidos educativos:} ampliar la sección \textit{Aprender} con materiales adicionales, ejemplos comparativos y recorridos temáticos.
    \item \textbf{Analítica educativa más profunda:} ofrecer métricas que ayuden a interpretar no solo el resultado financiero, sino también los patrones de decisión del usuario.
    \item \textbf{Nuevos escenarios de simulación:} extender el conjunto de activos, mercados o situaciones históricas disponibles para enriquecer el valor formativo del simulador.
    \item \textbf{Exportación e informes más completos:} mejorar la calidad y variedad de las salidas exportables, tanto para el historial como para el informe final del modo carrera.
\end{itemize}

Estas líneas de trabajo futuro se consideran razonables porque no duplican funcionalidades ya implementadas, sino que amplían capacidades reales sobre una base ya consolidada.

\section{Conclusión personal}

La realización del proyecto ha permitido integrar competencias de desarrollo web, gestión de datos, despliegue en la nube, validación técnica y diseño de interacción en un único sistema. Más allá del resultado funcional, el proceso ha puesto de manifiesto la importancia de iterar, validar sobre entorno real y documentar adecuadamente las decisiones adoptadas.

Asimismo, el proyecto ha servido para comprobar que la calidad final de una aplicación no depende únicamente de su lógica interna, sino también de la coherencia de la interfaz, de la estabilidad del despliegue y de la capacidad de corregir incidencias reales detectadas durante la validación final.
